% Lámina apaisada en 3 columnas para la hoja esquemática de métodos de análisis sintáctico (Temas 4-6)
\newgeometry{margin=0.85cm}
\begin{landscape}
\thispagestyle{empty}
\begingroup
\fontsize{8.2pt}{9.6pt}\selectfont
\setlength{\parindent}{0pt}
\setlength{\parskip}{1.5pt}
\setlength{\columnsep}{0.4cm}
\setlength{\columnseprule}{0.35pt}
\setlength{\tabcolsep}{2.5pt}
\renewcommand{\arraystretch}{1.05}
\setlength{\abovedisplayskip}{3pt}
\setlength{\belowdisplayskip}{3pt}
\setlength{\abovedisplayshortskip}{2pt}
\setlength{\belowdisplayshortskip}{2pt}
\let\olditemize\itemize
\def\itemize{\olditemize\setlength{\itemsep}{0.5pt}\setlength{\parskip}{0pt}\setlength{\topsep}{1.5pt}\setlength{\partopsep}{0pt}}
\let\oldenumerate\enumerate
\def\enumerate{\oldenumerate\setlength{\itemsep}{0.5pt}\setlength{\parskip}{0pt}\setlength{\topsep}{1.5pt}\setlength{\partopsep}{0pt}}
\newcommand{\hsec}[1]{\par\vspace{4pt}\noindent{\fontsize{9.4pt}{11pt}\bfseries #1}\par\vspace{2pt}}
\makeatletter
\renewcommand{\subsection}{\@ifstar\hsec\hsec}
\makeatother

\begin{center}
{\large\bfseries Procesadores de Lenguajes · Hoja Esquemática de Métodos de Análisis Sintáctico}\\[1pt]
{\scriptsize Condiciones, conjuntos, algoritmos y resolución de conflictos de los Temas 4, 5 y 6.}
\end{center}
\vspace{-4pt}

\begin{multicols*}{3}

\subsection*{0. Mapa de métodos}
\begin{tabular}{@{}llp{4.3cm}@{}}
\toprule
\textbf{Método} & \textbf{Tipo} & \textbf{Qué se calcula}\\
\midrule
LL(1) & descendente & $\mathrm{INI}$, $\mathrm{SIG}$, tabla $M[A,a]$\\
Prec.\ operador & ascendente & $\mathrm{PRIMERO}$, $\mathrm{ULTIMO}$ (terminales); rel.\ entre terminales\\
Prec.\ simple & ascendente & $\mathrm{PRIMERO}^+$, $\mathrm{ULTIMO}^+$ (todos los símbolos)\\
LR(0) & ascendente & ítems LR(0); reducir en toda fila\\
SLR(1) & ascendente & ítems LR(0); reducir solo en $\mathrm{SIG}$\\
LR(1) & ascendente & ítems con símbolo de anticipación\\
LALR(1) & ascendente & LR(1) con núcleos fusionados\\
\bottomrule
\end{tabular}

\smallskip
\textbf{No confundir:}
\begin{itemize}
\item $\mathrm{INI}/\mathrm{SIG}$: solo terminales (y $\$$ en $\mathrm{SIG}$); se usan en LL(1), SLR y LR(1).
\item $\mathrm{PRIMERO}/\mathrm{ULTIMO}$: primer/último \textbf{terminal}, admitiendo un no terminal delante/detrás. Solo precedencia de operador.
\item $\mathrm{PRIMERO}^+/\mathrm{ULTIMO}^+$: primer/último \textbf{símbolo cualquiera} ($V_N\cup V_T$). Solo precedencia simple.
\end{itemize}

\subsection*{1. Iniciales y seguidores (Tema 4)}
\[
\mathrm{INI}(\alpha)=\{a\in V_T\mid\alpha\Rightarrow^*a\beta\}\cup\{\varepsilon\text{ si }\alpha\Rightarrow^*\varepsilon\}
\]
\[
\mathrm{SIG}(A)=\{a\in V_T\cup\{\$\}\mid S\Rightarrow^*\gamma Aa\delta\}
\]
\textbf{Cálculo de $\mathrm{INI}$} (repetir hasta punto fijo):
\begin{enumerate}
\item $a$ terminal: $\mathrm{INI}(a)=\{a\}$.
\item $X\to\varepsilon$: $\varepsilon\in\mathrm{INI}(X)$.
\item $X\to Y_1Y_2\cdots Y_k$: añadir $\mathrm{INI}(Y_1)-\{\varepsilon\}$; si $Y_1\Rightarrow^*\varepsilon$, añadir $\mathrm{INI}(Y_2)-\{\varepsilon\}$, y así sucesivamente. Si todos son anulables, añadir $\varepsilon$.
\end{enumerate}
\textbf{Cálculo de $\mathrm{SIG}$} (repetir hasta punto fijo):
\begin{enumerate}
\item $\$\in\mathrm{SIG}(S)$ (axioma).
\item $A\to\alpha B\beta$: $\mathrm{INI}(\beta)-\{\varepsilon\}\subseteq\mathrm{SIG}(B)$.
\item $A\to\alpha B$, o $A\to\alpha B\beta$ con $\varepsilon\in\mathrm{INI}(\beta)$: $\mathrm{SIG}(A)\subseteq\mathrm{SIG}(B)$.
\end{enumerate}
\emph{Nunca} entra $\varepsilon$ en $\mathrm{SIG}$. Hay que revisar \emph{todas} las apariciones de $B$ en las partes derechas.

\columnbreak

\subsection*{2. Análisis LL(1) (Tema 4)}
\textbf{Requisitos previos:} sin recursividad por la izquierda (directa ni indirecta) y factorizada por la izquierda.
\begin{itemize}
\item Rec.\ izq.: $A\to A\alpha\mid\beta\;\Rightarrow\;A\to\beta A'$, $A'\to\alpha A'\mid\varepsilon$.
\item Factorización: $A\to\alpha\beta_1\mid\alpha\beta_2\;\Rightarrow\;A\to\alpha A'$, $A'\to\beta_1\mid\beta_2$.
\end{itemize}
\textbf{Tabla $M$:} para cada producción $A\to\alpha$:
\begin{itemize}
\item $M[A,a]=A\to\alpha$ para todo $a\in\mathrm{INI}(\alpha)$;
\item si $\varepsilon\in\mathrm{INI}(\alpha)$: $M[A,b]=A\to\alpha$ para todo $b\in\mathrm{SIG}(A)$.
\end{itemize}
\textbf{Condición LL(1):} ninguna casilla con dos producciones. Para $A\to\alpha\mid\beta$:
\[
\mathrm{INI}(\alpha)\cap\mathrm{INI}(\beta)=\emptyset
\]
y si $\beta\Rightarrow^*\varepsilon$, además $\mathrm{INI}(\alpha)\cap\mathrm{SIG}(A)=\emptyset$.
\textbf{Pila y traza:} empezar con $\$S$. Cima terminal igual a la entrada: emparejar. Cima no terminal $X$: sustituir por la parte derecha de $M[X,a]$ en orden inverso. Casilla vacía: error.

\subsection*{3. Desplazamiento-reducción (Tema 5)}
Acciones del analizador: \textbf{desplazar} el siguiente token a la pila, \textbf{reducir} sustituyendo el pivote por la parte izquierda, \textbf{aceptar} y \textbf{error}.\\
Conflictos posibles: \textbf{d/r} (desplazar frente a reducir) y \textbf{r/r} (duda entre dos reducciones).

\subsection*{4. Precedencia de operador (Floyd, Tema 5)}
\textbf{Condiciones:} (1) sin dos no terminales consecutivos en ninguna parte derecha; (2) sin $\varepsilon$-producciones; (3) relaciones disjuntas.

\textbf{Conjuntos} ($X\in V_N$, $a\in V_T$):
\[
\mathrm{PRIM}(A)=\{a\mid A\Rightarrow^+a\cdots\text{ o }A\Rightarrow^+Xa\cdots\}
\]
\[
\mathrm{ULT}(A)=\{a\mid A\Rightarrow^+\cdots a\text{ o }A\Rightarrow^+\cdots aX\}
\]
\textbf{Relaciones} (solo entre terminales y $\$$):
\begin{itemize}
\item $a\doteq b$ si $A\to\alpha a\beta b\gamma$ con $\beta\in V_N$ o $\beta=\varepsilon$.
\item $a\lessdot b$ si $A\to\alpha aB\beta$ y $b\in\mathrm{PRIM}(B)$.
\item $a\gtrdot b$ si $A\to\alpha Bb\beta$ y $a\in\mathrm{ULT}(B)$.
\item $\$\lessdot b$ para $b\in\mathrm{PRIM}(S)$; $a\gtrdot\$$ para $a\in\mathrm{ULT}(S)$.
\end{itemize}
\textbf{Algoritmo} ($a$ = terminal más alto en pila, $b$ = entrada):
\begin{itemize}
\item $a\lessdot b$ o $a\doteq b$: desplazar $b$.
\item $a\gtrdot b$: reducir sacando hasta que el nuevo terminal en cima cumpla $\lessdot$ respecto al último extraído.
\end{itemize}
\textbf{Limitaciones:} no admite doble precedencia (menos unario frente a binario) y puede aceptar cadenas erróneas (\texttt{id++id}) al no validar la parte derecha al reducir.

\columnbreak

\subsection*{5. Precedencia simple (Wirth--Weber, Tema 5)}
\textbf{Condiciones:} (1) sin dos producciones con la misma parte derecha; (2) sin $\varepsilon$-producciones; (3) relaciones disjuntas en $V\times V$.

\textbf{Conjuntos} ($Y\in V_N\cup V_T$):
\[
\mathrm{PRIM}^+(A)=\{Y\mid A\Rightarrow^+Y\cdots\}
\]
\[
\mathrm{ULT}^+(A)=\{Y\mid A\Rightarrow^+\cdots Y\}
\]
\textbf{Relaciones} para cada par adyacente $XY$ en una parte derecha:
\begin{itemize}
\item $X\doteq Y$.
\item si $Y\in V_N$: $X\lessdot Z$ para todo $Z\in\mathrm{PRIM}^+(Y)$.
\item si $X\in V_N$: $Z\gtrdot b$ para todo $Z\in\mathrm{ULT}^+(X)$ y todo terminal $b\in\{Y\}\cup\mathrm{PRIM}^+(Y)$.
\item $\$\lessdot Z$ para $Z\in\mathrm{PRIM}^+(S)$; $Z\gtrdot\$$ para $Z\in\mathrm{ULT}^+(S)$.
\end{itemize}
\textbf{Pivote:} subcadena delimitada entre $\lessdot$ y el primer $\gtrdot$ posterior. Se reduce por la producción cuya parte derecha coincide exactamente.\\
\textbf{Errores:} de precedencia (casilla vacía) y de reducción (el pivote aislado no coincide con ninguna parte derecha).

\smallskip
\begin{tabular}{@{}lp{2.85cm}p{3.45cm}@{}}
\toprule
 & \textbf{Operador} & \textbf{Simple}\\
\midrule
Tabla & terminales y $\$$ & $V_N\cup V_T$ y $\$$\\
Conjuntos & $\mathrm{PRIM}$, $\mathrm{ULT}$ & $\mathrm{PRIM}^+$, $\mathrm{ULT}^+$\\
Prohíbe & $\varepsilon$; dos $V_N$ seguidos & $\varepsilon$; partes der.\ iguales\\
Reducción & ignora parte derecha & pivote = parte der.\ exacta\\
\bottomrule
\end{tabular}

\end{multicols*}

\newpage
\thispagestyle{empty}

\begin{center}
{\large\bfseries Procesadores de Lenguajes · Hoja Esquemática (Análisis LR y YACC)}\\[1pt]
{\scriptsize Construcción de colecciones, fusión LALR(1), resolución de conflictos y recuperación de errores.}
\end{center}
\vspace{-4pt}

\begin{multicols*}{3}

\subsection*{6. Análisis LR: núcleo común (Tema 6)}
\textbf{Pila:} $s_0X_1s_1\cdots X_ms_m$. Con estado $s$ en cima y entrada $a$, se consulta $\mathrm{Acci\acute{o}n}[s,a]$:
\begin{itemize}
\item $d_j$: apilar $a$ y $j$; avanzar entrada.
\item $r_k$ ($A\to\beta$): desapilar $2|\beta|$ elementos; con estado $t$ debajo, apilar $A$ y $\mathrm{Ir\_a}[t,A]$.
\item $\mathrm{acc}$: aceptar. Casilla vacía: error.
\end{itemize}
\textbf{Gramática aumentada:} $(0)\; S'\to S$.\\
\textbf{Ítem LR(0):} $A\to\alpha\mathbin{\boldsymbol{\cdot}}\beta$. Si $\beta=\varepsilon$, es ítem de reducción ($A\to\varepsilon$ da $A\to\mathbin{\boldsymbol{\cdot}}$).\\
\textbf{Clausura:} si $A\to\alpha\mathbin{\boldsymbol{\cdot}} B\beta\in I$ y $B\to\gamma$, añadir $B\to\mathbin{\boldsymbol{\cdot}}\gamma$ hasta punto fijo.\\
\textbf{Ir\_a (Goto):} $\mathrm{goto}(I,X)=\mathrm{clausura}(\{A\to\alpha X\mathbin{\boldsymbol{\cdot}}\beta\mid A\to\alpha\mathbin{\boldsymbol{\cdot}} X\beta\in I\})$.\\
\textbf{Relleno común de la tabla:}
\begin{itemize}
\item $\mathrm{goto}(I_i,a)=I_j$, $a\in V_T\;\Rightarrow\;\mathrm{Acci\acute{o}n}[i,a]=d_j$.
\item $\mathrm{goto}(I_i,A)=I_j$, $A\in V_N\;\Rightarrow\;\mathrm{Ir\_a}[i,A]=j$.
\item $S'\to S\mathbin{\boldsymbol{\cdot}}\in I_i\;\Rightarrow\;\mathrm{Acci\acute{o}n}[i,\$]=\mathrm{acc}$.
\end{itemize}

\subsection*{7. LR(0) y SLR(1) (Tema 6)}
\begin{tabular}{@{}lp{7.1cm}@{}}
\toprule
\textbf{LR(0)} & $A\to\alpha\mathbin{\boldsymbol{\cdot}}\in I_i$: $r_k$ en \textbf{todas} las columnas de Acción.\\
\textbf{SLR(1)} & $A\to\alpha\mathbin{\boldsymbol{\cdot}}\in I_i$, $A\ne S'$: $r_k$ en $\mathrm{Acci\acute{o}n}[i,a]$ $\forall a\in\mathrm{SIG}(A)$.\\
\bottomrule
\end{tabular}

\textbf{Conflictos en SLR(1)} dentro de un estado $I$:
\begin{itemize}
\item \textbf{d/r:} $A\to\alpha\mathbin{\boldsymbol{\cdot}} a\beta\in I$ y $B\to\gamma\mathbin{\boldsymbol{\cdot}}\in I$ con $a\in\mathrm{SIG}(B)$.
\item \textbf{r/r:} $A\to\alpha\mathbin{\boldsymbol{\cdot}}\in I$ y $B\to\beta\mathbin{\boldsymbol{\cdot}}\in I$ con $\mathrm{SIG}(A)\cap\mathrm{SIG}(B)\ne\emptyset$.
\end{itemize}
En \textbf{LR(0)} estricto, todo estado con un ítem de reducción debe contener únicamente ese ítem.

\columnbreak

\subsection*{8. LR(1) canónico (Tema 6)}
\textbf{Ítem:} $[A\to\alpha\mathbin{\boldsymbol{\cdot}}\beta,\,a]$ con $a\in V_T\cup\{\$\}$.\\
$I_0=\mathrm{clausura}(\{[S'\to\mathbin{\boldsymbol{\cdot}} S,\,\$]\})$.

\textbf{Clausura LR(1):} si $[A\to\alpha\mathbin{\boldsymbol{\cdot}} B\beta,\,a]\in I$ y $B\to\eta$, añadir $[B\to\mathbin{\boldsymbol{\cdot}}\eta,\,b]$ para todo $b\in\mathrm{INI}(\beta a)$:
\begin{itemize}
\item $\beta$ no anulable: $b\in\mathrm{INI}(\beta)$ (no pasa $a$).
\item $\beta=\varepsilon$: $b=a$ (la anticipación se hereda).
\item $\beta$ anulable: $b\in(\mathrm{INI}(\beta)-\{\varepsilon\})\cup\{a\}$.
\end{itemize}
\textbf{Ir\_a (Goto):} avanza el punto y \textbf{copia} las anticipaciones.\\
Dos estados coinciden solo si son iguales tanto los ítems como sus anticipaciones.\\
\textbf{Reducción:} $[A\to\alpha\mathbin{\boldsymbol{\cdot}},\,a]\in I_i$ ($A\ne S'$): $\mathrm{Acci\acute{o}n}[i,a]=r_k$.

\subsection*{9. LALR(1) (Tema 6)}
\begin{enumerate}
\item Construir la colección LR(1).
\item Fusionar los estados con idéntico \textbf{núcleo} (mismos ítems LR(0)) uniendo sus conjuntos de anticipación ($I_{i,j}$).
\item Actualizar destinos de $\mathrm{goto}$ sobre los estados fusionados.
\end{enumerate}
\textbf{Teorema:} la fusión LALR nunca crea conflictos d/r nuevos, pero \textbf{puede introducir conflictos r/r}. Número de estados = número de estados de SLR.

\subsection*{10. Jerarquía y comparación LR}
\begin{tabular}{@{}lllp{1.3cm}@{}}
\toprule
 & \textbf{Ítems} & \textbf{Col.\ de $r_k$} & \textbf{Estados}\\
\midrule
LR(0) & sin anticip. & todas & $n$\\
SLR & sin anticip. & $\mathrm{SIG}(A)$ & $n$\\
LALR & fusionados & unión de ant. & $n$\\
LR(1) & con anticip. & solo $a$ & $\ge n$\\
\bottomrule
\end{tabular}

\smallskip
$\mathrm{LR}(0)\subset\mathrm{SLR}(1)\subset\mathrm{LALR}(1)\subset\mathrm{LR}(1)$, y $\mathrm{LL}(1)\subset\mathrm{LR}(1)$.\\
\emph{Ojo:} $\mathrm{LL}(1)$ \textbf{no} está contenida en $\mathrm{SLR}(1)$ ni en $\mathrm{LALR}(1)$. Ninguna gramática ambigua es $\mathrm{LL}(k)$ ni $\mathrm{LR}(k)$.

\columnbreak

\subsection*{11. Resolución de conflictos y YACC}
Estado con $[E\to E\,\theta_1\,E\mathbin{\boldsymbol{\cdot}}]$ y entrada $\theta_2$:

\smallskip
\begin{tabular}{@{}ll@{}}
\toprule
\textbf{Situación} & \textbf{Acción}\\
\midrule
$\theta_2$ tiene mayor precedencia & desplazar\\
$\theta_2$ tiene menor precedencia & reducir\\
igual precedencia, asoc.\ izquierda & reducir\\
igual precedencia, asoc.\ derecha & desplazar\\
\texttt{\%nonassoc} & error\\
\bottomrule
\end{tabular}

\smallskip
\textbf{Dangling else:} ante d/r con \texttt{else}, se elige \textbf{desplazar} para asociar cada \texttt{else} con el \texttt{if} más cercano.\\
\textbf{Conflictos r/r:} suelen deberse a producciones redundantes; se eliminan las reglas equivalentes.\\
\textbf{YACC:} genera analizadores LALR(1). Por defecto resuelve \textbf{d/r $\to$ desplazar} y \textbf{r/r $\to$ primera regla}. Las directivas \texttt{\%left}, \texttt{\%right}, \texttt{\%nonassoc} dan mayor precedencia cuanto más abajo aparecen; \texttt{\%prec} asigna precedencia explícita a una producción (por ejemplo, menos unario).

\subsection*{12. Recuperación de errores}
\textbf{LL(1):} modo pánico usando $\mathrm{SIG}(A)$ como conjunto de sincronización (sacar no terminal en \texttt{sinc}, descartar entrada en casilla vacía).\\
\textbf{LR:} rutinas específicas en casillas vacías; método de Graham--Rhodes (condensación del contexto y corrección).\\
\textbf{YACC:} regla \texttt{A : error w}. Saca estados hasta encontrar uno con $A\to\mathbin{\boldsymbol{\cdot}}\,\texttt{error}\,w$, desplaza \texttt{error}, descarta entrada hasta el token de sincronización $w$ (o un símbolo de $\mathrm{SIG}(A)$ si $w=\varepsilon$) y reduce a $A$. \texttt{yyerrok} reactiva la emisión inmediata de errores; \texttt{yyclearin} descarta el lookahead actual.

\end{multicols*}
\endgroup
\end{landscape}
\restoregeometry
